% Topological Casimir Effect
\documentclass[aps,prl,twocolumn,superscriptaddress,nofootinbib]{revtex4-2}

\usepackage{amsmath,amssymb,amsfonts}
\usepackage{hyperref}
\usepackage{booktabs}

\newcommand{\Spec}{\mathrm{Spec}}

\begin{document}

\title{Topological Protection of Vacuum Energy Sign:\\
A Testable Prediction from K-Theory\\
for Dirichlet--Neumann Casimir Scaling}

\author{Wright Brothers Research Program}
\affiliation{Independent Research, 2026}

\date{\today}

\begin{abstract}
The Casimir energy of a one-dimensional cavity changes sign when the
boundary conditions are changed from Dirichlet--Dirichlet (DD) to
Dirichlet--Neumann (DN): $E_{DD}\propto\zeta(-1)=-1/12<0$ (attractive)
versus $E_{DN}\propto\zeta_{\neg 2}(-1)=+1/12>0$ (repulsive).  This is
well established.  We observe that the localization
$\mathbb{Z}\to\mathbb{Z}[1/2]$ underlying the DN configuration
induces a change in algebraic K-theory,
$K_1(\mathbb{Z})=\mathbb{Z}/2\to K_1(\mathbb{Z}[1/2])=
\mathbb{Z}/2\times\mathbb{Z}$, introducing a topological
winding number that protects the sign of the vacuum energy.
If this protection is physical, the standard $1/a^4$ decay of
the Casimir energy density must exhibit a \emph{deviation} at a
critical separation $a_c$ where the topological transition occurs:
the density should plateau rather than continue to decay.  We
propose a minimal experiment---precision measurement of the
DN Casimir force scaling law---that can confirm or rule out this
prediction.  The required apparatus (Si-MEMS or circuit QED) exists.
\end{abstract}

\maketitle

% ===================================================================
\section{The Sign Flip}
% ===================================================================

Consider a massless scalar field in a one-dimensional cavity of
length~$a$.  With Dirichlet conditions at both ends (DD, ``$\lambda/2$''),
the mode frequencies are $\omega_n=n\pi c/a$, $n=1,2,3,\ldots$,
and the regularized vacuum energy is
\begin{equation}\label{eq:EDD}
  E_{DD}=\frac{\pi\hbar c}{2a}\,\zeta(-1)
  =-\frac{\pi\hbar c}{24a}<0.
\end{equation}
With Dirichlet at one end and Neumann at the other (DN,
``$\lambda/4$''), the modes are $\omega_n=(2n{-}1)\pi c/(2a)$,
$n=1,2,3,\ldots$, giving
\begin{equation}\label{eq:EDN}
  E_{DN}=\frac{\pi\hbar c}{4a}\sum_{n=1}^{\infty}(2n{-}1)
  =\frac{\pi\hbar c}{4a}\,\zeta_{\neg 2}(-1)
  =+\frac{\pi\hbar c}{48a}>0.
\end{equation}
Here $\zeta_{\neg 2}(s)\equiv(1{-}2^{-s})\zeta(s)
=\sum_{n\;\mathrm{odd}}n^{-s}$ is the Riemann zeta function with
the $p{=}2$ Euler factor removed.  The \emph{sign flip}
$E_{DD}<0\to E_{DN}>0$ is a standard result, confirmed by multiple
regularization methods~\cite{Boyer1974,Kenneth2002,Hertzberg2005,Milton2004}.

In three spatial dimensions, the analogous quantities involve
$\zeta(-3)=1/120$ and $\zeta_{\neg 2}(-3)=-7/120$, giving
Casimir energy densities
\begin{equation}
  \rho_{DD}=-\frac{\pi^2\hbar c}{720\,a^4},\qquad
  \rho_{DN}=+\frac{7\pi^2\hbar c}{720\,a^4}.
\end{equation}

% ===================================================================
\section{The K-Theoretic Observation}
% ===================================================================

The passage from DD to DN corresponds algebraically to the
localization $\mathbb{Z}\to\mathbb{Z}[1/2]$: one removes all
integers divisible by~2 from the mode sum, retaining only the
odd modes.  This localization induces a change in algebraic
K-theory~\cite{Quillen1972}:
\begin{equation}\label{eq:K1}
  K_1(\mathbb{Z})=\mathbb{Z}/2
  \;\longrightarrow\;
  K_1(\mathbb{Z}[1/2])=\mathbb{Z}/2\times\mathbb{Z}.
\end{equation}
The emergent $\mathbb{Z}$ factor is a \emph{winding number}.
In condensed-matter physics, an analogous winding number
(the SSH invariant~\cite{SSH1979}) topologically protects
edge states: a system with winding number $w\neq 0$ cannot
be continuously deformed to $w=0$ without closing the gap.

We conjecture that the $\mathbb{Z}$ winding number
in~\eqref{eq:K1} \emph{topologically protects the sign} of the
vacuum energy in the DN configuration: once the system is in the
$E>0$ (repulsive) phase, it cannot be continuously deformed back
to $E<0$ (attractive) without a discrete topological transition.

% ===================================================================
\section{The Prediction: Scaling Deviation}
% ===================================================================

In standard quantum field theory, the Casimir energy density
$\rho\propto 1/a^4$ for both DD and DN boundary conditions.
As the plate separation $a$ increases, $\rho$ decays to zero.
The sign difference between DD and DN is maintained at every~$a$,
but the \emph{magnitude} vanishes as $a\to\infty$.

If the sign of the vacuum energy is topologically protected, we
predict a qualitatively different behavior: there exists a
critical separation $a_c$ such that for $a>a_c$, the energy
density $\rho_{DN}$ \emph{ceases to decay} as $1/a^4$ and instead
plateaus at a constant value determined by the topological
invariant.

\begin{center}
\begin{tabular}{lcc}
\toprule
Regime & Standard QED & Topological \\
\midrule
$a<a_c$ & $\rho\propto 1/a^4$ & $\rho\propto 1/a^4$ \\
$a>a_c$ & $\rho\propto 1/a^4$ & $\rho\to\mathrm{const}$ \\
\midrule
Scaling exponent & $-4$ everywhere & $-4\to 0$ at $a_c$ \\
\bottomrule
\end{tabular}
\end{center}

The critical separation $a_c$ is not predicted by our analysis.
It could range from nanometers (if set by the lattice of
$\Spec(\mathbb{Z})$) to macroscopic scales (if set by
cosmological parameters).  The prediction is the \emph{existence}
of the deviation, not its location.

% ===================================================================
\section{Analogy: Bulk--Edge Correspondence}
% ===================================================================

In topological insulators, the \emph{bulk--edge correspondence}
guarantees that a material with nontrivial bulk topological
invariant must have gapless edge states, regardless of surface
details.  The edge states are ``protected'' by the bulk.

We propose an analogous \emph{bulk--boundary correspondence} for
the vacuum: the sign of the Casimir energy (a boundary effect) is
protected by the K-theoretic invariant of the bulk vacuum
($K_1(\mathbb{Z}[1/2])$).  Just as the conductance of a
topological edge state is quantized and insensitive to disorder,
the sign of $E_{DN}$ should be insensitive to the details of the
boundary, depending only on the topological class.

This would manifest experimentally as a deviation from the
$1/a^4$ scaling: the ``boundary'' effect (Casimir energy) would
acquire a ``bulk'' contribution (constant energy density) once
the topological protection is established.

% ===================================================================
\section{Experimental Test}
% ===================================================================

\noindent\textbf{Minimal experiment:} Measure the DN Casimir force
as a function of plate separation $a$ over at least two decades
(e.g., $a=50$\,nm to $5\,\mu$m) and fit the power law
$F\propto a^{-\alpha}$.

\begin{itemize}
\item Standard QED predicts $\alpha=5$ (force $\propto -dE/da
\propto a^{-5}$ in 3D) for all~$a$.
\item The topological prediction is $\alpha<5$ for $a>a_c$,
ultimately approaching $\alpha=1$ (constant density $\to$
force $\propto 1/a$... actually $F=-d(\rho V)/da$ if $\rho$
constant then $F=-\rho A$ independent of $a$, i.e.,
$\alpha\to 0$).
\end{itemize}

Any confirmed deviation of $\alpha$ from~5 in the DN
configuration, not present in the DD configuration, would be
evidence for topological protection of the vacuum energy sign.

\medskip
\noindent\textbf{Platforms:}
\begin{itemize}
\item \textbf{Si-MEMS} (Princeton-type~\cite{Princeton2017}):
Geometric Casimir force measurement at 50--500\,nm.  DN boundary
conditions achievable via T-shaped or open-ended structures.
\item \textbf{Circuit QED}: $\lambda/4$ vs.\ $\lambda/2$
superconducting resonator with tunable end impedance $Z$.
Sweep $Z$ from 0 (Dirichlet) to $\infty$ (Neumann) and measure
Lamb shift scaling.
\end{itemize}

% ===================================================================
\section{Discussion}
% ===================================================================

We emphasize that this paper makes a \emph{single, falsifiable
prediction}: the DN Casimir force deviates from $1/a^5$ scaling
at some separation $a_c$.  If no deviation is found across the
accessible range, the topological protection conjecture is
excluded for $a_c$ within that range.

The prediction is independent of the broader ``$\Spec(\mathbb{Z})$
as spacetime'' program.  It relies only on:
(i)~the algebraic K-theory computation~\eqref{eq:K1}, which is a
theorem;
(ii)~the physical identification of $K_1$ winding with vacuum
energy sign, which is the conjecture.

If confirmed, the implications extend far beyond the Casimir
effect: a topologically protected constant vacuum energy density
would constitute a \emph{controllable dark energy}, with
applications to spacetime engineering.

\begin{thebibliography}{99}
\bibitem{Boyer1974} T.~H.~Boyer, Phys.\ Rev.\ A \textbf{9}, 2078 (1974).
\bibitem{Kenneth2002} O.~Kenneth \textit{et al.}, Phys.\ Rev.\ Lett.\ \textbf{89}, 033001 (2002).
\bibitem{Hertzberg2005} M.~P.~Hertzberg \textit{et al.}, Phys.\ Rev.\ Lett.\ \textbf{95}, 250402 (2005).
\bibitem{Milton2004} K.~A.~Milton, J.\ Phys.\ A \textbf{37}, R209 (2004).
\bibitem{Quillen1972} D.~Quillen, in \textit{Algebraic K-Theory I}, Springer LNM \textbf{341} (1973).
\bibitem{SSH1979} W.~P.~Su, J.~R.~Schrieffer, A.~J.~Heeger, Phys.\ Rev.\ Lett.\ \textbf{42}, 1698 (1979).
\bibitem{Princeton2017} L.~Tang \textit{et al.}, Nature Photon.\ \textbf{11}, 97 (2017).
\end{thebibliography}

\end{document}
